\textbf{Proof.}
We divide the argument into the orbit reduction, the parity classification, the unequal two-block calculation, and the homogeneous calculation.

Let \(\mathfrak G_{\mathcal P}\) be the product of the within-block permutation groups.  For \(A\in\mathcal C_n\), define
\[
\overline A=|\mathfrak G_{\mathcal P}|^{-1}
\sum_{g\in\mathfrak G_{\mathcal P}}U_gAU_g^\top.
\tag{1}
\]
The covariance hull and every score-quality shell for \(S_{\mathcal P}\) are invariant under this action.  Since \(A\mapsto q_\rho(A;S_{\mathcal P})\) is a maximum of linear functionals, it is convex, and Jensen's inequality gives
\[
q_\rho(\overline A;S_{\mathcal P})\leq q_\rho(A;S_{\mathcal P}).
\tag{2}
\]

For an assignment \(z\), write \(r_j=\sum_{i\in G_j}z_i\).  Conditional on \(r\), uniform within-block permutation gives
\[
\mathbb E[z_i z_{i'}\mid r]=\frac{r_jr_\ell}{m_jm_\ell}
\quad
(i\in G_j,\ i'\in G_\ell,\ j\ne\ell),
\tag{3}
\]
and, for distinct \(i,i'\in G_j\),
\[
\mathbb E[z_i z_{i'}\mid r]
=\frac{r_j^2-m_j}{m_j(m_j-1)}.
\tag{4}
\]
Equation \((4)\) follows by expanding \(r_j^2\).  Equations \((3)\)--\((4)\) are exactly the entries of \(A_{\mathcal P}(rr^\top)\).  Thus a convex representation of \(A\) supplies \(G\in\mathcal G_{\mathcal P}\) with \(\overline A=A_{\mathcal P}(G)\), proving the full-hull converse by \((2)\).

Conversely, express \(G\) as a mixture of \(rr^\top\) over \(\mathcal R_{\mathcal P}\).  For every such \(r\), choose signs with those block sums, permute them uniformly within blocks, and multiply the whole assignment by an independent fair sign.  This is a mean-zero law in \(\mathcal D_n\), and \((3)\)--\((4)\) show that its covariance is \(A_{\mathcal P}(rr^\top)\).  Mixing realizes \(A_{\mathcal P}(G)\).  Compactness proves attainment and hence
\[
v^*(\rho;S_{\mathcal P})
=\min_{G\in\mathcal G_{\mathcal P}}
q_\rho(A_{\mathcal P}(G);S_{\mathcal P}).
\tag{5}
\]

For parity, if every \(m_j\) is even, the lattice contains \(r=0\); its covariance annihilates all block-constant centered vectors, so \(\kappa(S_{\mathcal P})=0\).  Conversely, if a cut-hull mixture has zero compression on \(S_{\mathcal P}\), positivity implies that every supported assignment satisfies \(x^\top z=0\) for every \(x\in S_{\mathcal P}\).  Its block-sum vector \(r\) is therefore orthogonal to every \(t\) satisfying \(\sum_jm_jt_j=0\), hence is proportional to \((m_1,\ldots,m_B)\).  Balance gives \(\sum_jr_j=0\), so the proportionality constant is zero.  Thus every block sum is zero, which is possible only when every block size is even.  This proves the parity equivalence.

We now solve the unequal two-block problem.  Write the block sizes as \(m<\ell\).  Balance forces \(r=(R,-R)\), where
\[
R\in\{-m,-m+2,\ldots,m\}.
\]
Because \(m+\ell=n\) is even, the block sizes have the same parity.  The convex hull of the possible squares is
\[
q=\mathbb E[R^2]\in[e,m^2],
\qquad
e=0\ \text{in the even--even case},\quad e=1\ \text{in the odd--odd case}.
\tag{6}
\]
Every point of this interval is obtained by mixing squared lattice values.

Let \(w\) be the unit vector in the one-dimensional score space, constant on each block.  Direct normalization gives the score eigenvalue
\[
a(q)=w^\top A_{\mathcal P}(G)w=\frac{nq}{m\ell}.
\tag{7}
\]
On the within-block contrast space of a block of size \(j>1\), the eigenvalue is
\[
b_j(q)=\frac{j^2-q}{j(j-1)}.
\tag{8}
\]
These mutually orthogonal spaces decompose \(H_n\).  If \(m>1\), put \(B(q)=\max\{b_m(q),b_\ell(q)\}\).  For \(s=\rho^2\), allocating the nonscore energy to the larger contrast eigenvalue gives
\[
q_{\sqrt{s}}(A_{\mathcal P}(q);S_{\mathcal P})
=\max\{a(q),\,sa(q)+(1-s)B(q)\}.
\tag{9}
\]

The three eigenvalues meet at
\[
q_c=\frac{m\ell}{n-1},
\qquad
a(q_c)=b_m(q_c)=b_\ell(q_c)=c_n.
\tag{10}
\]
This is also the complete-randomization moment, because summing the covariance \(c_nP_H\) over either block gives \(m\ell/(n-1)\).  Subtracting the affine expressions in \((7)\)--\((8)\) shows that, for \(e\leq q\leq q_c\),
\[
b_m(q)\geq b_\ell(q)\geq a(q),
\]
whereas \(a(q)\) is largest for \(q\geq q_c\).  Consequently the minimization of \((9)\) reduces on \([e,q_c]\) to
\[
sa(q)+(1-s)b_m(q).
\tag{11}
\]
Its slope changes sign at
\[
s_0=\frac{\ell}{m(n-1)}.
\tag{12}
\]
Thus \((11)\) is minimized at \(q_c\) for \(s\leq s_0\), and at \(q=e\) for \(s\geq s_0\).  Substitution gives
\[
v^*(\sqrt{s};S_{\mathcal P})
=\min\{c_n,(1-s)b_0+s\kappa_2\},
\quad
b_0=\frac{m^2-e}{m(m-1)},
\quad
\kappa_2=\frac{ne}{m\ell}.
\tag{13}
\]
The two branches meet at \(s_0\).  The parity-endpoint laws in the statement realize \(q=e\), while complete randomization realizes \(q=q_c\).  Equation \((13)\) gives \(\kappa(S_{\mathcal P})=\kappa_2\), and \(\text{thm:admissibility-frontier}\) gives \(\gamma(S_{\mathcal P})=s_0\) and \(\rho^*(S_{\mathcal P})=\sqrt{s_0}\).

For fixed \(q\), the loss in \((9)\) is convex in \(s\), and the oracle value in \((13)\) is concave in \(s\).  Their difference is therefore convex, so its maximum on \([0,1]\) occurs at an endpoint.  Since \(v(0)=c_n\) and \(v(1)=\kappa_2\),
\[
\Delta_{\mathrm{reg}}(A_{\mathcal P}(q);S_{\mathcal P})
=\max\{\max(a(q),B(q))-c_n,\ a(q)-\kappa_2\}.
\tag{14}
\]
On \([e,q_c]\), the first term is \(b_m(q)-c_n\), which decreases, and the second is \(a(q)-\kappa_2\), which increases.  They are equal at
\[
q_{\mathrm{reg}}=(1-s_0)e+s_0q_c,
\tag{15}
\]
with common value
\[
s_0(c_n-\kappa_2).
\tag{16}
\]
Outside \([e,q_c]\), equation \((14)\) is no smaller.  The orbit reduction makes this a converse for every covariance in \(\mathcal C_n\).  Mixing the parity-endpoint law with complete randomization using weight \(s_0\) gives \((15)\), so the converse is attained.  Scaling \((16)\) by \(4M/n\) proves the stated regret formula and shows that the constraints at \(\rho=0\) and \(\rho=1\) are jointly binding.

If \(m=1\), equation \((6)\) forces \(q=1=q_c\).  Only the \(\ell\)-block contrast space remains, and \((7)\)--\((8)\) both equal \(c_n\).  Hence the loss is \(c_n\) at every quality, \(\kappa=c_n\), \(\gamma=1\), \(\rho^*=1\), and the regret is zero.

Finally, suppose there are \(B\) equal blocks of even size \(m=n/B\).  Averaging also over block permutations reduces a covariance to
\[
aP_S+b(P_H-P_S),
\qquad pa+hb=n,
\]
where \(p=B-1\) and \(h=n-B\).  Its score-shell loss is \(\max\{a,sa+(1-s)b\}\).  When \(a\leq c_n\), one has \(b\geq a\) and
\[
sa+(1-s)b
=\frac{n}{h}(1-s)+a\left\{s-\frac{p}{h}(1-s)\right\}.
\]
The coefficient of \(a\) changes sign at \(\beta=p/(n-1)\), whereas \(a>c_n\) already has loss above the complete-randomization value.  Complete randomization realizes \(a=b=c_n\), and independent within-block complete randomization realizes \(a=0\), \(b=\delta_m=n/h\).  Therefore
\[
v^*(\sqrt{s};S_B)=\min\{c_n,\delta_m(1-s)\},
\qquad
\kappa(S_B)=0,
\qquad
\gamma(S_B)=\beta,
\qquad
\rho^*(S_B)=\sqrt\beta.
\tag{17}
\]
The mixture with weight \(\beta\) on complete randomization has normalized regret \(c_n\beta=np/(n-1)^2\).  Conversely, for any \(A\in\mathcal C_n\), let \(K=\lambda_{\max}(P_SAP_S\mid S)\) and \(L=\lambda_{\max}(A\mid H_n)\).  If \(r=\Delta_{\mathrm{reg}}(A;S_B)\), the endpoint constraints give \(r\geq K\) and \(r\geq L-c_n\), while the trace gives \(n\leq pK+hL\).  Hence \(r\geq np/(n-1)^2\), proving the matching full-hull lower bound.  Multiplication by \(4M/n\) gives \(R^*(n,S_B)=4Mp/(n-1)^2\).

Setting \(B=n/2\) and block size two in \((17)\) gives the displayed pair specialization.  For unequal partitions with more than two blocks, equation \((5)\) remains the exact finite program, and no further closed form has been asserted. \(\square\)